\documentclass[10pt,
% a4paper,
%twocolumn,
fleqn,
%landscape, 
% papersize,
dvipdfmx,
uplatex
]{jsarticle}



\def\maru#1{\textcircled{\scriptsize#1}}%丸囲み番号

% \RequirePackage[2020/09/30]{platexrelease}

%太字設定
\usepackage[deluxe]{otf}

\usepackage{emathEy}

\usepackage[g]{esvect}

%大きな文字
\usepackage{fix-cm}

%定理環境
\usepackage{emathThm}
%\theoremstyle{boxed}
\theorembodyfont{\normalfont}
\newtheorem{Question}{問題}[subsection]
\newtheorem{Q}{}[subsection]
\newtheorem{question}[Question]{}
\newtheorem{quuestion}{}[subsection]

%セクション，大問番号のデザイン
\renewcommand{\labelenumi}{(\arabic{enumi})}
\renewcommand{\theenumii}{\alph{enumii})}
\renewcommand{\thesection}{第\arabic{section}章}

%用紙サイズの詳細設定
\usepackage{bxpapersize}
\papersizesetup{size={80mm,45mm}}
\usepackage[top=0.7zw,bottom=0truemm,left=3truemm,right=133truemm]{geometry}
\usepackage[dvipdfmx]{graphicx}

%余白など
\usepackage{setspace} % 行間
\setlength{\mathindent}{1zw}
\setlength\parindent{0pt}


%色カラーに関する設定
\usepackage{color}
\definecolor{shiro}{rgb}{0.95703125,0.87109375,0.7421875}
\definecolor{kin}{rgb}{0.95703125,0.87109375,0.7421875}
\definecolor{orange}{rgb}{1,0.7,0.2}
\definecolor{bradorange}{rgb}{1,0.5,0}
\definecolor{pink}{rgb}{0.9176,0.5686,0.5960}
\definecolor{mizu}{rgb}{0.6156,0.8,0.9955}
\color{kin}
% \pagecolor{hukamido}

\usepackage{at}%図の配置
% \usepackage{wallpaper}

\begin{document}



\at(0cm,0cm){\includegraphics[width=8cm,bb=0 0 1920 1080]{./youtube/thumbnails/templates/smart_background/図形と方程式.jpeg}}
{\color{orange}\bf\boldmath\huge\underline{線形計画法}}\vspace{0.3zw}

\LARGE 
\bf\boldmath 問．$3x+y\geqq 6,\;x+3y\geqq 6,\;x+y\leqq 6$のとき，$x+2y$の最大値と最小値を求めよ．
\at(6.6cm,0.2cm){\small\color{bradorange}$\overset{\text{図形と方程式}}{\text{典型}}$}


\newpage



\at(0cm,0cm){\includegraphics[width=8cm,bb=0 0 1920 1080]{./youtube/thumbnails/templates/smart_background/図形と方程式.jpeg}}
{\color{orange}\bf\boldmath\huge\underline{不等式の表す領域}}\vspace{0.3zw}

\small 
\bf\boldmath 問．次の不等式が表す領域を図示せよ．\\
(1)  $y>x^2+1$\\
(2)  $3x-2y-2\geqq 0$\\
(3)  $x\leqq 2$\\
(4)  $\left(x+2\right)^2+y^2<1$\\
(5)  $x^2+y^2-6x-2y+1\geqq 0$\\
(6)  $x^2+y^2<{25},\;y<3x-5$\\
(7)  $\left(x-y\right)\left(x+y-2\right)<0$\\

\at(6.6cm,0.2cm){\small\color{bradorange}$\overset{\text{図形と方程式}}{\text{典型}}$}


\newpage



\at(0cm,0cm){\includegraphics[width=8cm,bb=0 0 1920 1080]{./youtube/thumbnails/templates/smart_background/図形と方程式.jpeg}}
{\color{orange}\bf\boldmath\LARGE\underline{三角形の重心の軌跡}}\vspace{0.3zw}

\LARGE 
\bf\boldmath 問．$2$点$\text{O}\left(0,\;0\right),\;\text{A}\left(1,\;0\right)$と円$x^2+y^2=9$上を動く点$\text{Q}$を頂点とする三角形$\text{OAQ}$の重心$\text{P}$の軌跡を求めよ．
\at(6.6cm,0.2cm){\small\color{bradorange}$\overset{\text{図形と方程式}}{\text{典型}}$}


\newpage



\at(0cm,0cm){\includegraphics[width=8cm,bb=0 0 1920 1080]{./youtube/thumbnails/templates/smart_background/図形と方程式.jpeg}}
{\color{orange}\bf\boldmath\huge\underline{軌跡 Lv.1 }}\vspace{0.3zw}

\LARGE 
\bf\boldmath 問．$2$点$\text{A}\left(-3,\;0\right),\;\text{B}\left(2,\;0\right)$からの距離の比が$3:2$であるような点$\text{P}$の軌跡を求めよ．
\at(6.6cm,0.2cm){\small\color{bradorange}$\overset{\text{図形と方程式}}{\text{典型}}$}


\newpage



\at(0cm,0cm){\includegraphics[width=8cm,bb=0 0 1920 1080]{./youtube/thumbnails/templates/smart_background/図形と方程式.jpeg}}
{\color{orange}\bf\boldmath\LARGE\underline{円に内接$\cdot$外接する円}}\vspace{0.3zw}

\huge 
\bf\boldmath 問．中心が$\left(4,\;3\right)$で円$x^2+y^2=1$に接する円の方程式を求めよ．
\at(6.6cm,0.2cm){\small\color{bradorange}$\overset{\text{図形と方程式}}{\text{典型}}$}


\newpage



\at(0cm,0cm){\includegraphics[width=8cm,bb=0 0 1920 1080]{./youtube/thumbnails/templates/smart_background/図形と方程式.jpeg}}
{\color{orange}\bf\boldmath\LARGE\underline{円に引いた接線の方程式}}\vspace{0.3zw}

\large 
\bf\boldmath 問．次の接線の方程式を求めよ．\\
(1)  円$x^2+y^2=5$上の点$\left(3,\;4\right)$における接線\\
(2)  点$\left(1,\;3\right)$から円$x^2+y^2=5$に引いた接線\\

\at(6.6cm,0.2cm){\small\color{bradorange}$\overset{\text{図形と方程式}}{\text{典型}}$}


\newpage



\at(0cm,0cm){\includegraphics[width=8cm,bb=0 0 1920 1080]{./youtube/thumbnails/templates/smart_background/図形と方程式.jpeg}}
{\color{orange}\bf\boldmath\LARGE\underline{円と直線が接する条件}}\vspace{0.3zw}

\LARGE 
\bf\boldmath 問．円$x^2+y^2={10}$と直線$y=2x+m$が接するとき，定数$m$の値を求めよ．
\at(6.6cm,0.2cm){\small\color{bradorange}$\overset{\text{図形と方程式}}{\text{典型}}$}


\newpage



\at(0cm,0cm){\includegraphics[width=8cm,bb=0 0 1920 1080]{./youtube/thumbnails/templates/smart_background/図形と方程式.jpeg}}
{\color{orange}\bf\boldmath\Large\underline{円と直線が共有点をもつ条件}}\vspace{0.3zw}

\LARGE 
\bf\boldmath 問．円$x^2+y^2=8$と直線$y=x+m$が共有点を持つとき，定数$m$の値の範囲を求めよ．
\at(6.6cm,0.2cm){\small\color{bradorange}$\overset{\text{図形と方程式}}{\text{典型}}$}


\newpage



\at(0cm,0cm){\includegraphics[width=8cm,bb=0 0 1920 1080]{./youtube/thumbnails/templates/smart_background/図形と方程式.jpeg}}
{\color{orange}\bf\boldmath\huge\underline{円と直線の共有点}}\vspace{0.3zw}

\Large 
\bf\boldmath 問．円$x^2+y^2=5$と次の直線の共有点の座標を求めよ．\\
(1)  $y=x-1$\\
(2)  $y=2x+5$\\

\at(6.6cm,0.2cm){\small\color{bradorange}$\overset{\text{図形と方程式}}{\text{典型}}$}


\newpage



\at(0cm,0cm){\includegraphics[width=8cm,bb=0 0 1920 1080]{./youtube/thumbnails/templates/smart_background/図形と方程式.jpeg}}
{\color{orange}\bf\boldmath\LARGE\underline{$3$点を通る円の方程式}}\vspace{0.3zw}

\large 
\bf\boldmath 問．$3$点$\text{A}\left(2,\;1\right),\;\text{B}\left(6,\;3\right),\;\text{C}\left(-1,\;2\right)$がある．\\
(1)  $3$点$\text{A},\;B,\;C$を通る円の方程式を求めよ．\\
(2)  三角形$\text{ABC}$の外心の座標と，外接円の半径を求めよ．\\

\at(6.6cm,0.2cm){\small\color{bradorange}$\overset{\text{図形と方程式}}{\text{典型}}$}


\newpage



\at(0cm,0cm){\includegraphics[width=8cm,bb=0 0 1920 1080]{./youtube/thumbnails/templates/smart_background/図形と方程式.jpeg}}
{\color{orange}\bf\boldmath\LARGE\underline{$2$直線の平行$\cdot$垂直条件}}\vspace{0.3zw}

\Large 
\bf\boldmath 問．$2$直線$ax+2y=1,\;x+\left(a-1\right)y=3$が次の条件を満たすとき，定数$a$の値を求めよ．\\
(1)  平行\\
(2)  垂直\\

\at(6.6cm,0.2cm){\small\color{bradorange}$\overset{\text{図形と方程式}}{\text{典型}}$}


\end{document}

